\usepackage[provide=*,spanish]{babel}
\usepackage{fontspec}
\defaultfontfeatures{Ligatures=TeX}
\setmainfont{Latin Modern Roman}
\setsansfont{Latin Modern Sans}
\setmonofont{Latin Modern Mono}
\usepackage[a4paper,margin=2.25cm,headheight=15pt]{geometry}
\usepackage{amsmath,amssymb}
\usepackage{booktabs,longtable,tabularx,array}
\usepackage{graphicx,float}
\usepackage{xcolor}
\usepackage{hyperref}

\definecolor{uvgverde}{HTML}{176B52}
\definecolor{uvgoscuro}{HTML}{103F33}
\definecolor{grisclaro}{HTML}{F2F5F4}
\definecolor{azulcli}{HTML}{17324D}
\hypersetup{colorlinks=true,linkcolor=uvgoscuro,urlcolor=uvgverde,pdfauthor={Grupo 1 - CC3067},pdftitle={Laboratorio 3 - Algoritmos de Enrutamiento}}

\pagestyle{plain}
\renewcommand{\arraystretch}{1.18}

\newcolumntype{Y}{>{\raggedright\arraybackslash}X}
\newcolumntype{C}[1]{>{\centering\arraybackslash}p{#1}}
\newenvironment{calculo}[1][]{\par\medskip\noindent\color{uvgoscuro}\hrule\smallskip\textbf{Desarrollo del calculo}\par\color{black}}{\par\smallskip\color{uvgoscuro}\hrule\color{black}\medskip}
\newenvironment{decision}[1][]{\par\medskip\noindent\color{uvgoscuro}\hrule\smallskip\textbf{Decision de diseno}\par\color{black}}{\par\smallskip\color{uvgoscuro}\hrule\color{black}\medskip}
\newenvironment{evidencia}[1][]{\par\medskip\noindent\color{azulcli}\hrule\smallskip\textbf{Evidencia real de Packet Tracer}\par\color{black}}{\par\smallskip\color{azulcli}\hrule\color{black}\medskip}

\newcommand{\universidad}{Universidad del Valle de Guatemala}
\newcommand{\curso}{CC3067 Redes}
\newcommand{\seccion}{Seccion 10}
\newcommand{\grupo}{Grupo 1}
\newcommand{\integrantes}{%
\begin{tabular}{ll}
Pablo Daniel Barillas Moreno & 22193\\
Andres Rafael Chivalan & 21534
\end{tabular}}

\newcommand{\portadalab}[2]{%
\begin{titlepage}
\centering
{\Large\bfseries \universidad\par}
\vspace{1.2cm}
{\large \curso\ -- \seccion\par}
\vspace{2.2cm}
{\Huge\bfseries Laboratorio 3\par}
\vspace{0.35cm}
{\LARGE Algoritmos de Enrutamiento\par}
\vspace{0.8cm}
{\Large\color{uvgverde}#1\par}
\vfill
{\large\bfseries \grupo\par}
\vspace{0.45cm}
\integrantes
\vfill
{\large #2\par}
\end{titlepage}}

